%! Setting Up --- Learning from Data — Unit 8, Lesson 8.0 — Warm-Up
\documentclass[10pt]{article}
\usepackage{linalg-article}
\usepackage{linalg-boxes}
\newcommand{\TallMath}[1]{$\displaystyle #1\rule[-1.4em]{0pt}{3.2em}$}

\begin{document}

\pageheader{Unit 8, Lesson 8.0}{Warm-Up}
\namedateperiod

\vspace{0.12in}

\begin{notesbox}{Getting Ready: Ideas We'll Use Today}
\small These three ideas --- two from earlier units, one from plain arithmetic --- are the pieces a
neural network is built from. Answer quickly.

\vspace{4pt}
\textbf{1. Matrix times a vector (Unit 1).} Compute $A\mathbf{v}$ for
$A=\begin{bmatrix} 1 & -1 \\ 1 & 1 \end{bmatrix}$ and $\mathbf{v}=\begin{bmatrix} 1 \\ 2 \end{bmatrix}$:
\[
  A\mathbf{v}=\begin{bmatrix} 1 & -1 \\ 1 & 1 \end{bmatrix}\begin{bmatrix} 1 \\ 2 \end{bmatrix}
  = \begin{bmatrix}\rule{0pt}{1.1em}\blank{1.0cm}\\[2pt]\blank{1.0cm}\end{bmatrix}.
\]

\vspace{4pt}
\textbf{2. How far off? (Unit 4).} Least squares measures error by \emph{squaring} it. A guess of $4$
misses a target of $6$. The error is $4-6=\blank{1.2cm}$, and its \emph{square} is $\blank{1.2cm}$.
(Squaring makes the error positive and punishes big misses.)

\vspace{4pt}
\textbf{3. The max function (arithmetic).} Compute:
\[
  \max(0,\,5)=\blank{1.0cm},\qquad \max(0,\,-2)=\blank{1.0cm},\qquad \max(0,\,0)=\blank{1.0cm}.
\]
\end{notesbox}

\end{document}
